\begingroup
\renewcommand{\arraystretch}{1.45}
\Needspace{32\baselineskip}
\InterpolationArticle{35}
If the function $X$ does not stop at $\sin nx$, but continues further, with the following terms expressed by $\delta^{\mu+1}\sin(\mu+1)x+\delta^{\mu+2}\sin(\mu+2)x$ etc., we shall have:
\[
\begin{aligned}
&\zeta'={}\frac{1}{\sin\mu a}\{\delta'\sin\mu a+\delta^{\mu+1}\sin 2\mu a+\delta^{2\mu+1}\sin 3\mu a+\text{etc.}\\[5pt]
&\qquad+\delta^{2\mu-1}\sin\mu a+\delta^{3\mu-1}\sin 2\mu a+\delta^{4\mu-1}\sin 3\mu a+\text{etc.}\}\\[5pt]
&\zeta''={}\frac{1}{\sin\mu a}\{\delta''\sin\mu a+\delta^{\mu+2}\sin 2\mu a+\delta^{2\mu+2}\sin 3\mu a+\text{etc.}\\[5pt]
&\qquad+\delta^{2\mu-2}\sin\mu a+\delta^{3\mu-2}\sin 2\mu a+\delta^{4\mu-2}\sin 3\mu a+\text{etc.}\}\\[5pt]
&\zeta'''={}\frac{1}{\sin\mu a}\{\delta'''\sin\mu a+\delta^{\mu+3}\sin 2\mu a+\delta^{2\mu+3}\sin 3\mu a+\text{etc.}\\[5pt]
&\qquad+\delta^{2\mu-3}\sin\mu a+\delta^{3\mu-3}\sin 2\mu a+\delta^{4\mu-3}\sin 3\mu a+\text{etc.}\}
\end{aligned}
\]
etc. For the last coefficient, however,
\[
\begin{aligned}
&\zeta^{\mu}=\frac{1}{\sin\mu a}(\delta^{\mu}\sin\mu a+\delta^{2\mu}\sin 2\mu a+\delta^{3\mu}\sin 3\mu a+\text{etc.})
\end{aligned}
\]
\Needspace{40\baselineskip}
These formulas show how far the difference between $X$ and $X''$ may be neglected. This latter formula will be the simplest among all similar ones that satisfy the $\mu$ proposed values; all these, however, will be contained in the formula $X''+Y\sin x(\cos\mu x-\cos\mu a)$, where $Y$ denotes, as in Article 32, an undetermined function of the arc $x$ free of sines. And generally, if $X'''$ is any function of the same form as $X$, i.e. containing only sines, that satisfies the $\mu$ given values, the formula $X'''+Y\sin x(\cos\mu x-\cos\mu a)$ will contain all functions of this kind; if $Y$ is suitably determined, it can be reduced to the $\mu^{\mathrm{th}}$ order, and in this way the function $X''$ itself must necessarily result. In precisely the same way as in Article 32, the following general rule is obtained for this purpose: for any term in $X'''$ of the form $L\sin\lambda x$, where $\lambda$ is greater than $\mu$, one must substitute in $X''$, setting $\lambda=k\mu+\lambda'$ so that $k\mu$ is the multiple of $\mu$ immediately below $\lambda$, and thus $\lambda'$ lies between the limits $1$ and $\mu$ inclusive, the terms
\[
\begin{aligned}
&L\frac{\sin(k+1)\mu a}{\sin\mu a}\sin\lambda'x+L\frac{\sin k\mu a}{\sin\mu a}\sin(\mu-\lambda')x
\end{aligned}
\]
which, whenever $\lambda'=\mu$, reduce to the single term $L\frac{\sin\lambda a}{\sin\mu a}\sin\mu x$.
\SourceLeaf{327}
\Needspace{25\baselineskip}
\InterpolationArticle{36}
The transformations presented in the preceding article and in Article 32 can be derived more neatly from a certain general theorem; since it is also quite elegant in itself, we shall briefly touch on it here as well.

\textsc{Theorem.} \textit{Let $\lambda,\lambda',\lambda''$ denote any integers. Let $\mu$ denote an integer that divides the differences between them, $\lambda'-\lambda,\lambda''-\lambda',\lambda-\lambda''$ (e.g. unity), $x$ an undetermined arc, and $a$ a fixed arc: the functions}
\[
\begin{aligned}
&P=\sin(\lambda'-\lambda'')a\cos\lambda x+\sin(\lambda''-\lambda)a\cos\lambda'x+\sin(\lambda-\lambda')a\cos\lambda''x\\[5pt]
&Q=\sin(\lambda'-\lambda'')a\sin\lambda x+\sin(\lambda''-\lambda)a\sin\lambda'x+\sin(\lambda-\lambda')a\cos\lambda''x
\end{aligned}
\]
\textit{will be divisible by} $\cos\mu x-\cos\mu a$\textit{.}

\Needspace{29\baselineskip}
\textit{Proof.} When $\lambda$ is divisible by $\mu$, and thus $\lambda',\lambda''$ will also be divisible by $\mu$, it is easily confirmed that the value of $P,Q$, on substituting $x=a$, is identically $=0$; hence $P$ will not be changed by the following substitutions:
\[
\begin{array}{lll}
\cos\lambda x&\text{is replaced by}&\cos\lambda x-\cos\lambda a\\
\cos\lambda'x&&\cos\lambda'x-\cos\lambda'a\\
\cos\lambda''x&&\cos\lambda''x-\cos\lambda''a
\end{array}
\]
nor will $Q$ by the substitutions:
\[
\begin{array}{lll}
\sin\lambda x&\text{is replaced by}&\displaystyle\sin\lambda x-\frac{\sin\lambda a\sin\mu x}{\sin\mu a}\\[6pt]
\sin\lambda'x&&\displaystyle\sin\lambda'x-\frac{\sin\lambda'a\sin\mu x}{\sin\mu a}\\[6pt]
\sin\lambda''x&&\displaystyle\sin\lambda''x-\frac{\sin\lambda''a\sin\mu x}{\sin\mu a}
\end{array}
\]
\Needspace{32\baselineskip}
But these six expressions are divisible by $\cos\mu x-\cos\mu a$; it suffices to show this for the first and fourth for a positive value of $\lambda$. Indeed, it is easily confirmed by multiplication that
\[
\begin{aligned}
&\sin\mu a(\cos\lambda x-\cos\lambda a)\\[5pt]
&=(\cos\mu x-\cos\mu a)\{2\sin\mu a\cos(\lambda-\mu)x+2\sin 2\mu a\cos(\lambda-2\mu)x\\[5pt]
&\quad \qquad+2\sin 3\mu a\cos(\lambda-3\mu)x+\text{etc.}+2\sin(\lambda-\mu)a\cos\mu x+\sin\lambda a\}
\end{aligned}
\]
\[
\begin{aligned}
&\sin\mu a\sin\lambda x-\sin\lambda a\sin\mu x\\[5pt]
&=(\cos\mu x-\cos\mu a)\{2\sin\mu a\sin(\lambda-\mu)x+2\sin 2\mu a\sin(\lambda-2\mu)x\\[5pt]
&\quad \qquad+2\sin 3\mu a\sin(\lambda-3\mu)x+\text{etc.}+2\sin(\lambda-\mu)a\sin\mu x\}
\end{aligned}
\]
The case where $\lambda$ is negative reduces of itself to this one. Hence it is clear that the functions
\SourceLeaf{328}
$P$, $Q$ are composed of parts divisible by $\cos\mu x-\cos\mu a$, and are therefore themselves also divisible by this divisor.

\Needspace{34\baselineskip}
II. When $\lambda$ is not divisible by $\mu$, let $l$ be an arbitrary integer divisible by $\mu$, and set $\lambda=l+\theta$, $\lambda'=l'+\theta$, $\lambda''=l''+\theta$; then $l'$, $l''$ will also be divisible by $\mu$. Now it is clear that, on setting
\[
\begin{aligned}
&\sin(l'-l'')a\cos lx+\sin(l''-l)a\cos l'x+\sin(l-l')a\cos l''x=P'\\[5pt]
&\sin(l'-l'')a\sin lx+\sin(l''-l)a\sin l'x+\sin(l-l')a\sin l''x=Q'
\end{aligned}
\]
there results
\[
\begin{aligned}
&P=P'\cos\theta x-\theta'\sin\theta x,\quad Q=P'\sin\theta x+Q'\cos\theta x,
\end{aligned}
\]
and that the functions $P'$, $Q'$, since they belong to the first case already completed, are divisible by $\cos\mu x-\cos\mu a$: hence evidently $P$ and $Q$ will also be divisible by $\cos\mu x-\cos\mu a$. Q.~E.~D. However, the proof of the first case has been carried out in such a way that it cannot be applied without some explanations whenever $\sin\mu a=0$; but then $P=0$, $Q=0$, so that no proof at all is needed.

\Needspace{24\baselineskip}
Thus, if $\lambda-\lambda'=k\mu$, $\lambda-\lambda''=(k+1)\mu$ are set, it is clear that the following are divisible by $\cos\mu x-\cos\mu a$:
\[
\begin{aligned}
&\sin\mu a\cos\lambda x-\sin(k+1)\mu a\cos\lambda'x+\sin k\mu a\cos(\lambda'-\mu)x\\[5pt]
&\sin\mu a\sin\lambda x-\sin(k+1)\mu a\sin\lambda'x+\sin k\mu a\sin(\lambda'-\mu)x
\end{aligned}
\]
and hence also
\[
\begin{aligned}
&\displaystyle\cos\lambda x-\frac{\sin(k+1)\mu a\cos\lambda'x-\sin k\mu a\cos(\mu-\lambda')x}{\sin\mu a}\\[5pt]
&\displaystyle\sin\lambda x-\frac{\sin(k+1)\mu a\sin\lambda'x+\sin k\mu a\sin(\mu-\lambda')x}{\sin\mu a}
\end{aligned}
\]
From this the reason for the substitutions in Articles 32, 35 becomes clear at once: for the quotient from the latter division will consist only of sines, and hence will evidently again be divisible by $\sin x$.

\endgroup
